\chapter{Introduction}
\label{chap:introduction}

\section{Background}
\label{section:background}

Communication is an essential part of daily life, yet people with hearing disabilities often face significant 
challenges when interacting with others. Many members of the deaf community rely on sign language as 
their primary form of communication,but most hearing people do not understand it.
This gap creates communication barriers in situations such as education, healthcare, employment, and emergency services.

In Thailand, many individuals with hearing disabilities depend on Thai Sign Language,
but professional interpreters are limited and not always available, especially in urgent situations.
As a result, communication between individuals with hearing disabilities and hearing users can be difficult.

Recent advances in artificial intelligence and machine learning provide new opportunities to address this problem.
Technologies such as speech recognition, and natural language processing enable systems to interpret gestures and speech automatically.

This project proposes TermTem, an AI-based communication assistant that integrates sign language recognition, speech recognition, 
conversational AI, and real time animation displayed through avatar to support real time communication between hearing disability individuals and hearing users.

\section{Problem Statement}
\label{section:problem-statement}

People with hearing disabilities often experience communication difficulties when interacting with individuals 
who do not understand Thai Sign Language. In many real life situations such as hospitals, workplaces, public services, 
or emergency cases, communication becomes difficult because sign language interpreters are not always available.
Existing solutions mainly rely on human interpreters or limited translation tools, which may be costly, unavailable, 
or not suitable for real time communication. 

Therefore, there is a need for an intelligent system that can assist in communication. Such a system should be able to translate
sign language, speech, and text in real time to reduce communication barriers and improve accessibility for people with hearing disabilities.

\section{Solution Overview}
\label{section:solution-overview}

This project proposes TermTem, an AI-based communication assistant that integrates sign language recognition, speech recognition, 
conversational AI, and real time animation displayed through avatar to support real time communication between hearing disability individuals and hearing users.

\subsection{Features}
\label{subsection:features}

\begin{enumerate}[leftmargin=80pt]
    \item Sign Language Recognition: The system can recognize Thai sign language gestures from video input and convert them into text.
    \item Speech Recognition: Speech from hearing users can be converted into text so that the system can process and interpret the message.
    \item Sign Language Animation Generation: The system can generate sign language animations through an avatar to help users with hearing disabilities understand spoken language.
    \item Conversational AI Assistant: The system integrates a conversational AI model that can understand user inputs and generate meaningful responses.
    \item Real-time Communication Support: The system is designed to provide near real time responses to enable smooth communication between users.
\end{enumerate}

\section{Target User}
\label{section:target-user}

People with hearing disabilities: 

\quad Demographics: Individuals of all ages in Thailand who primarily use Thai Sign Language as their main form of communication.

\quad Skill level: Basic smartphone proficiency.

\quad Industry or Domain: Daily communication scenarios such as healthcare services, public services, workplaces, and emergency situations.

Learners of Thai sign language: 

\quad Demographics: Students, educators, family members, or individuals interested in learning Thai Sign Language.

\quad Skill level: Beginners in sign language who require clear visual demonstrations and accurate sign language animations for learning.

\quad Industry or Domain: Education.
\newpage
\section{Terminology}
\label{section:terminology}
\begin{table}[h]
\centering
\renewcommand{\arraystretch}{1.5} % Adds padding for readability
\begin{tabular}{|l|p{10cm}|}
\hline
\textbf{Term} & \textbf{Description} \\ \hline
Thai Sign Language  & The primary sign language used by the deaf community in Thailand. \\ \hline
Sign Language Recognition & A computer vision technique used to detect and interpret sign language gestures. \\ \hline
Speech Recognition & Technology that converts spoken language into text. \\ \hline
Text-to-Speech  & Technology that converts text into spoken audio. \\ \hline
Conversational AI & AI technology that enables systems to understand and respond to human language. \\ \hline
Avatar Animation & A digital character used to visually represent sign language gestures. \\ \hline
\end{tabular}
\end{table}